\documentclass[a4paper,twocolumn]{article}
%\documentclass[10pt, conference, compsocconf]{IEEEtran}
%\documentclass[a4paper, conference]{IEEEtran}


\usepackage{pslatex} % -- times instead of computer modern, especially for the plain article class
\usepackage[colorlinks=false,bookmarks=false]{hyperref}
\usepackage{booktabs}
\usepackage{graphicx}
\usepackage{xcolor}
\usepackage{multirow}
\usepackage{cite}
%\usepackage{flushend} % even out the last page, but use only at the end when there is a bibliography

\newcommand{\code}[1]{{\small{\texttt{#1}}}}

% fatter TT font
\renewcommand*\ttdefault{txtt}
% another TT, suggested by Alex
% \usepackage{inconsolata}
% \usepackage[T1]{fontenc} % needed as well?

\usepackage{listings}

\newcommand{\todo}[1]{{\emph{TODO: #1}}}
\newcommand{\martin}[1]{{\color{blue} Martin: #1}}
\newcommand{\abcdef}[1]{{\color{red} Author2: #1}}

% uncomment following for final submission
%\renewcommand{\todo}[1]{}
%\renewcommand{\martin}[1]{}
%\renewcommand{\author2}[1]{}

\begin{document}

\title{Hardware Design with Chisel}

\author{Martin and maybe Jonatan}

% Most conferences have their own commands for author headings.

%\author{\IEEEauthorblockN{Edgar Lakis, Martin Schoeberl}\\
%\IEEEauthorblockA{Department of Applied Mathematics and Computer Science\\
%Technical University of Denmark\\
%Email: \texttt{edgar.lakis@gmail.com}, \texttt{masca@imm.dtu.dk}}
%}


\maketitle \thispagestyle{empty}

\begin{abstract}
A good paper has a short, concise abstract. The abstract
states the field of the research, the purpose, and the findings (results).
\end{abstract}


\section{Introduction}
\label{sec:intro}

\todo{A brief introduction what the paper is about. It shall include briefly the
main contributions and findings. The contributions can be bullet listed.}

This paper... \todo{purpose statement, latest in 4th paragraph}

\todo{Test the bib with a reference that gives background on time-predictable
computer architecture~\cite{tpca:jes}.}


\begin{figure}
  \centering
  \includegraphics{figures/broken-mux}

  \caption{An inkscape figure.}
  \label{fig:mux}
\end{figure}

\begin{figure}
  \centering
  \includegraphics{figures/mux}

  \caption{A figure from graffle.}
  \label{fig:mux}
\end{figure}

A paper is cited \cite{paper:example}.



\section{Conclusion}
\label{sec:conclusion}

\todo{Rephrase what this paper is about and list the main contributions and results.}

\subsection*{Acknowledgment}

\todo{Sometimes we received some help. Sometimes external funding.}



\subsection*{Source Access}

\martin{I love doing papers with available source under an
open-source license. It gives credit and good karma.}

\bibliographystyle{plain}
% Please do not add any references to msbib.bib.
% They get lost when I 'generate' is again (see Makefile)
\bibliography{chisel,msbib}

\end{document}
